\label{sec:conclusion}

M.2 operationalizes physical land use as a community character signal through
the USGS NLCD 2021 land cover classification.
The four-component land use vector---residential, commercial, open developed,
and natural---captures the physical landscape character of each census tract.
Cosine similarity between adjacent tracts modulates edge weights so that
METIS cuts preferentially at physical land use transitions.
The hard-cut rule for majority-water adjacencies encodes the well-established
legal doctrine that water bodies are natural community boundaries.

Together, these design choices address a gap in algorithmic redistricting:
the widespread practice of weighting edges by boundary length, with no
regard for the physical character of the communities on either side.
A long shared boundary between a residential neighborhood and an adjacent
commercial corridor is a natural cut seam, not a community bond.
M.2 makes this distinction explicit and quantitative.

\paragraph{Relationship to the M-track series.}
M.2 addresses the physical layer of community character: what is actually
built on the ground.
It complements M.1 (economic character: what jobs are located there),
M.3 (housing character: what kind of housing stock and households occupy
it), and M.6 (administrative zone membership: what service districts
govern it).
These four signals together cover the primary community dimensions that
residents identify when they describe why their neighborhood is a
recognizable community.

The M.8 composite index \citep{m8composite} assigns M.2 a default weight
of 0.20 (tied with M.1 as the second-highest weight after M.6).
This reflects two judgments: land use is a strong community signal
(weight 0.20) but slightly weaker than zone co-membership in legal
proceedings (weight 0.30), and at least as strong as economic character
signals that are better supported by employment data (M.1 at 0.20).

\paragraph{The hard-cut rule as a model for legal defensibility.}
The water hard-cut rule in M.2 is a model for how legal doctrine can be
encoded as an algorithmic constraint.
Rather than treating water boundaries as a soft preference, the rule
encodes the legal reality: cross-water districts are presumptively invalid
and require specific geographic justification (a bridge, a tunnel, an
established ferry service with contiguous community use).
The zero-weight encoding means the algorithmic output either satisfies
the water boundary constraint or is infeasible; there is no gradation
between compliance and violation.

Future M-track papers may adopt similar hard-cut rules for other legally
established boundaries: county line preservation requirements (where
applicable), school district preservation criteria (M.6), and
VRA-required minimum threshold adjacencies.
The pattern is: when a criterion is binary in legal practice
(you may not split this boundary), encode it as a binary constraint
in the algorithm, not as a soft weight.

\paragraph{Phase 2 commitments.}
Phase~2 will complete the following quantitative results, all currently
marked ${}^{\dagger}$:
\begin{enumerate}
  \item Land-use boundary coincidence (LUB) metric for NC-14, WI-8,
    and TX-38, comparing land-use weights to geographic baseline.
  \item Hard-cut validation: zero cross-water districts in NC (Pamlico Sound)
    and TX (Gulf Coast bays) under land-use weights.
  \item Polsby-Popper compactness comparison between land-use and geographic
    weight modes.
  \item Multi-seed variance analysis for all three states (Phase~3 if
    single-seed Phase~2 results are confirmatory).
\end{enumerate}

Phase~2 depends on NLCD 2021 GeoTIFF download and zonal statistics computation.
Once the NLCD data is available in the bisect data cache, the full pipeline
run can be completed in a single \texttt{bisect build} invocation with
\texttt{--weights-override land-use}.
